\section{Polymaths and Self-Creation}

\begin{quotex}
A way stands open into heaven, but none can enter the way except those who have heaven within them.
\flright{\textsc{Emanuel Swedenborg}}
\end{quotex}


It is very difficult to become an expert in one field. To become expert in two or more fields is extraordinary. Such a person is called a polymath and several polymaths have always intrigued me. The most obvious recent examples are Leonardo da Vinci\footnote{\url{https://en.wikipedia.org/wiki/Leonardo\_da\_Vinci}} and Johann Wolfgang Goethe. Da Vinci's accomplishments in art, science, and technology are well known. Goethe was a dramatist and poet, while his scientific studies in fields like optics and biology have been neglected. A book recommended by \textbf{Wolfgang Smith}, \emph{The Wholeness of Nature} by \textbf{Henri Bortoft}, provides a good explanation of Goethe's science of conscious participation in nature. Bortoft's training in visualization provided the key to understand Goethe.

Some will use a related term, the \textbf{Renaissance Man}, in place of polymath, although that is more like a well-rounded man, who does well in many disparate areas of his life, rather than as a specialist in a few. \textbf{Baldesar Castiglione}, in \emph{The Book of the Courtier}, describes this type, as we summarized in The Duty of the Wise Man\footnote{\url{http://www.gornahoor.net/?p=3187}}. Castiglione combines the duties of the courtier with the virtuous life. The goal is to consciously create a ``single whole'' out of one's life, which in most men is fragmented, random, and factitious. There are some polymaths whose lives I have studied in some depth at different points in my life.

\textbf{Albert Schweitzer} was a medical doctor who practiced as a medical missionary in Africa, much of which he documented. He made original contributions in theology, in philosophy through his ideal of Reverence for Life, and in the music of J S Bach. He was also an expert in church organs as well as a proficient instrumentalist.

\textbf{Alfred North Whitehead} was an expert in logic, math, science, and philosophy. After a career as a mathematics professor, he switched to philosophy. He opposed the Cartesian theory of bifurcation, and integrated qualities and creativity into a scientific worldview. At that time, I was more interested in holism and evolutionistic views. Nevertheless, just as \textbf{Arthur Lovejoy} showed how the idea of the Great Chain of Being was flattened out into an evolutionary worldview in time, it may be possible to make Whitehead's metaphysics more hierarchical. You can see how this might look in our posts on how Providence works in relationship to Will and Destiny. At every moment, Providence offers us creative ``possibilities of manifestation'' (Guenon's term), which we can integrate into our life. This does not require Whitehead's notion of a God in evolution.

\paragraph{Swedenborg}


Since we've been discussing the role of ``imagination'' in our Monday Gnosis seminars, this is opportune to bring up another polymath, \textbf{Emmanuel Swedenborg}. He began as a scientist and engineer, making contributions to metallurgy, geometry, chemistry, physiology, and then became interested in philosophy. Late in life, he began to have visionary and clairvoyant experiences which he described in several books.

His writings were important among Martinist and Hermetic groups, although \textbf{Valentin Tomberg} makes just a passing reference to him. Yet his ideas resonated beyond that. Curiously, the Zen Buddhist \textbf{D T Suzuki} wrote about Swedenborg, translated his books into Japanese, and referred to him as the ``Buddha of the North''. I suppose this connection to Buddhism should not be surprising, as Swedenborg referred to some lost teachings among the Sages of Tibet and of Tartary. I learned this from \textbf{Rene Guenon}\footnote{\url{http://www.gornahoor.net/?p=972}}. This idea was confirmed by the Catholic mystic \textbf{Anne Catherine Emmerich}. In the \textbf{Mountain of the Prophet}, also in Tibet, she saw the holy books of all the Traditions revealed to mankind. The point is that ``Mankind was not yet capable of receiving these gifts, another must first come.'' That led me to explore Tibetan Buddhism, as well as the writings of \textbf{Alice Bailey} and \textbf{Nicholas Roerich} to find those secret teachings. Perhaps I did not look hard enough and stopped too soon. It is an unfinished project, which I've hinted at, to use Madhyamika philosophy to express the metaphysics of Tradition.

\textbf{Henry Corbin}, the scholar of esoteric Islam, was also intensely interested in Swedenborg's ideas. He demonstrated a correspondence between Swedenborg's visions and mystical visions described in Islam. For them, there are three worlds: the sensory world, the world of the imagination, and the world of the intellect. The intermediate world of the imagination is the one described by the seers. A full review of those writings might be a good topic for the future. I know the interest will not be widespread, but it will be very intense for a few of us.

Nevertheless, the American Sufi Charles Upton, \emph{pace} Corbin, writes this about Swedenborg:

\begin{quotex}
That a great deal of profound doctrines can be found in Swedenborg's writings is undeniable. But now that the scriptures and classic of the world religions and the writings of history's greatest sages are readily available, we no longer need to take him as uniquely inspired authorities since we can judge them in the light of their orthodox originals ... there is no longer any reason to rely upon suspect sources. ... Since such doctrines were not available to a Swedish Lutheran of the 18\textsuperscript{th} century, it was necessary to reintroduce them via direct inspiration.
\end{quotex}


While it is true that we don't need to take Swedenborg as ``uniquely inspired'', it is nevertheless still the case that his contributions to our understanding of the imaginal or astral realm has not been reintroduced into the mainstream of Western Tradition. \textbf{E F Schumacher} mentioned the visionaries \textbf{Edgar Cayce} and \textbf{Jakob Lorber} in \emph{The Guide for the Perplexed}, as guides to his fields of knowledge.

\paragraph{Self-Creation}

When I was a young man, I realized my life as a mathematician was becoming detrimental to my social life. Time series differential equations or obscure issues in algebraic topology are conversation stoppers. So I resolved to become a latter day courtier, a Renaissance Man of sorts. I already has some expertise in mathematics and physics, I was more literary than most, and fortunately I was reasonably athletic. Although, as a proper Bostonian, I was taught ballroom dancing, I had to learn the Hustle, and later, Salsa, on my own.

In order to take on certain qualities I thought I was missing in my life, I began an intense study of the works of \textbf{Harold Robbins} -- certainly not great literature, but a good description of a world that was alien to me. But I always took literature personally. Earlier, I had read most of \textbf{Jack Kerouac}'s works and then \textbf{Henry Miller} and \textbf{Lawrence Durrell}; I adopted what I read into my own life and self-image. I saw myself in Joseph Conrad's books. From Hesse, I decided I was a combination of \textbf{Steppenwolf} and \textbf{Siddhartha}; my life course has been close to both.


This process took several years. I developed a solid workout program, I took creative writing classes, a Dale Carnegie course, seminars at the Florida Speaker's Association, and even acting classes. I was motivated by \textbf{Nietzsche}'s injunction:

\begin{quotex}
To ``give style'' to one's character-- a great and rare art! It is practiced by those who survey all the strengths and weaknesses of their nature and then fit them into an artistic plan until every one of them appears as art and reason and even weaknesses delight the eye. Here a large mass of second nature has been added; there a piece of original nature has been removed --- both times through long practice and daily work at it. Here the ugly that could not be removed is concealed; there it has been reinterpreted and made sublime. Much that is vague and resisted shaping has been saved and exploited for distant views; it is meant to beckon toward the far and immeasurable. In the end, when the work is finished, it becomes evident how the constraint of a single taste governed and formed everything large and small. Whether this taste was good or bad is less important than one might suppose, if only it was a single taste!
\end{quotex}


A human birth is too precious to let one's life be governed solely by chance. I learned the doctrinal basis for my task only later on. Of course, I had done an experiment recommended by Aleister Crowley as a way of freeing up one's personality; this was described in Aghori of the Mind\footnote{\url{http://www.gornahoor.net/?p=3214}}.

Tibetan Buddhism involves visualization practices. A common one is to see oneself in a mandala of the Buddha, gods, goddesses, and enlightened beings. I asked my guru what the purpose of that was. He explained that we are constantly creating a self-image subconsciously, without thought. We can create a new self-image as an enlightened being.

In high school they used to draw a curtain across the gymnasium to separate the boys from the girls in physical education classes. I used to visualize a different education during those classes. But that sort of daydreaming is passive and a waste of our sexual energies. Visualizing myself as a sort of Renaissance Man, a ``man about town'', successful in business and with women, did bring me a not insignificant degree of self-mastery, as I became able to overcome an inherent introversion in situations where it is not suitable.

Ultimately, then, I could visualize myself in an icon of the communion of saints like the Buddhists do. Beyond that, there is the world of pure intellect, beyond all images and thoughts.


\flrightit{Posted on 2016-05-30 by Cologero}

\begin{center}* * *\end{center}

\begin{footnotesize}
\begin{sffamily}
\texttt{Mr Man on 2016-05-30 at 04:08 said:}

Lookin good, Toni!

\hfill

\texttt{Tom on 2016-05-30 at 11:30 said:}

Fascinating insight into your personal development. I too used to read Kerouac and Miller and am guilty of styling myself as a Milleresque figure in my early twenties. I think perhaps it is hard not to take literature personally at that stage in life when one is so ready to be influenced. I wonder, perhaps self indulgently, how I might have benefited from earlier exposure to more spiritual literature and Traditional perspectives but fate goes ever as it must.

\hfill

\texttt{Alistair Fraser on 2016-05-30 at 18:53 said:}

'' I had to learn the Hustle, and later, Salsa, on my own. `

Looking more like the ``american gigolo'' in that image cologero
... ... . its good for the soul :-)


Thx for the rundown of some previous applications -- a lot of pro-active movements there , a possible review of the said correspondence sounds of importance

\hfill

\texttt{francismercuri on 2016-05-30 at 22:17 said:}

Interesting piece. Something I take away from it is a description of a movement from life stage to life stage:

Beginning with ``polymath'', man engages in ``doing something'' in the world, and (attempting/striving to) become a master of it--so, the ``perfection of the human state''.

Further beyond worldly and skilled ``prowess'' comes the need to clean up the mind and imagination; mastery of the ``astral'', correcting sloppy and slipshod use of the apparently never ending stream of consciousness we all take for granted--akin to what Guenon describes as ``psychic regeneration'', or as is represented in certain Initiatic cycles as a ``second degree'', ``fellow'', or ``journeyman''.

Lastly, one arrives at the point where all acquired worldly skill, and visualization expertise conclude, perhaps surprisingly--at ``nothing''--from ``successful polymath'' to an ``accomplished'' apophatic.

In tandem with this post's personal nature/disclosure, I think a very useful spiritual exercise involves keeping a diary. A regular diary is a surefire tool in the game of developing self-awareness.

\hfill

\texttt{O on 2016-12-09 at 19:51 said:}

``Perhaps I did not look hard enough and stopped too soon. It is an unfinished project, which I've hinted at, to use Madhyamika philosophy to express the metaphysics of Tradition.''

If this is an unfinished project, why do you discourage those of us with a firm grounding in the Western Traditional metaphysicists as well as a strong attraction to Buddhism from pursuing such a course in terms of initiation and practical traditional affiliation? For what you are mentioning won't happen if no one coming from our position delves deeply enough into those Vajrayana mysteries, which is only possible if studying and practicing under the guidance of a qualified master.

In terms of this particular goal, you certainly stopped too soon. The Tibetan tradition is ultra-complex, with diverse branches and many dimensions. While the Gelug order is good for dry scholastic Madhyamaka, it is not the best place to look for the most profound yogic mysteries. The Zhentong Madhyamaka Tantra approach of the more obscure Jonangpas is also much closer to the viewpoint of Western Traditionalists.

\hfill

\texttt{O on 2016-12-09 at 20:15 said:}

Suzuki's praise for Swedenborg as ``the Buddha of the North'' is curious to say the least. He was certainly a `genius' (a questionable concept as Schuon demonstrated in `To have a Centre') in his way, and a man of unusual psychic powers, but that doesn't equal true `enlightenment'/awakening or gnosis. Doctrinally he was hardly orthodox. To Fr Seraphim Rose, he was a kind of counter-initiate. I think it's safe to say that his overall influence on the spiritual landscape of Europe has been subversive; which isn't to say that a man of good discernment cannot find valid element in his works.

\hfill

\texttt{Cologero on 2016-12-11 at 10:37 said:}

Bear in mind that the original impulse came from Murti's ``Central Philosophy of Buddhism'', in which he relates Madhyamkia to Western thought. Murti was also knowledgeable about Guenon and referred explicitly to the idea of using Madhyamkia rather than the Vedanta. This is a very early post here:
\textit{Buddhism and Tradition}\footnote{\url{http://www.gornahoor.net/?p=299}}


This was written five years later: \textit{The Esoteric Path}\footnote{\url{http://www.gornahoor.net/?p=7872}}, and it may explain something to you.

As a matter of practicality, I cannot work on a couple of blogs, do translations, etc., and also work on this project. I doubt I would discourage anyone, also I would encourage him to at least start with Murti's book, even if it seems dry and scholastic.

From a personal point of view, I've been disappointed with the Dalai Lama who seems to have become Westernized and liberalized. Perhaps that should not hinder others.

As for actually changing ``affiliation'', I would discourage that simply on traditional grounds, and furthermore it is not even necessary to embark on this project. For example, the Medieval scholastics did not change their affiliation to ``pagan'' in order to embrace Aristotle and Neo-platonism.

I'll concede that, in this disordered age, the situation may not apply and a man has to do what is best for himself. However, in a traditional hierarchical society, men don't question their own tradition. It just makes no sense, and to say otherwise is just the result of the lingering effects of liberalism. As Evola wrote: in order to command, one must obey. So, again on a personal (and traditional) basis, I embrace the Nordic-Roman Tradition of the Middle Ages.

\hfill

\texttt{Cologero on 2016-12-12 at 12:49 said:}

Re Guenon on Swedenborg, please see \textit{Finding Agarttha}\footnote{\url{http://www.gornahoor.net/?p=972}}. What is the ``lost word'' that apparently eluded Swedenborg?

\hfill

\texttt{O on 2016-12-22 at 14:36 said:}

Thank you, Cologero, that clarifies it. I had somehow got the idea that you wanted someone thoroughly schooled in the metaphysics of the Mahayana to take on Tradition as a whole in the universal esoteric manner of a Guénon or Schuon. No such representative exists today as far as I am aware, although there have been previous examples, notably Marco Pallis, who was based in the Vajrayana, but whose literary production was rather limited.

Now, to adapt and organically integrate the Madhyamaka view into the specifically Western, Christian tradition is an endeavour whose degree of difficulty vastly exceeds that other project, but the zhentong view is much more suited for this than Madhyamaka as formulated by the Gelug mainstream, which easily comes across as rather nihilistic, although it isn't if subtly understood. I honestly doubt that I would have the required genius for such a daring and pioneering endeavour. I can also foresee that it would be met with massive scepticism from traditional Christians.

For example, could this task hypothetically be performed by someone in the priesthood? The thought has occasionally struck me whether becoming an Orthodox priest\footnote{There happens to be an institution here that makes such a course very straight forward... \url{http://www.sanktignatios.org/home-f20r3}} would have suited me, as being of a contemplative nature, that is the formal traditional caste in Europe I can most readily sympathise with. (I would never consider the Catholic priesthood due to a number of problems, ranging from the Pope and Vatican II to the fact that they cannot be married.) I can't help but wonder if being in such an official position would place too many limits on what I may do in terms of visible esoteric endeavours. What do you think? Could you have been open with your own perspectives, for example, as a priest?

\hfill

\texttt{O on 2016-12-22 at 14:58 said:}

I forgot to add that, in regards to the Dalai Lama, I share your dissatisfaction. His rather naive attitude to the West, modernity and profane science should be given an antidote through introducing him to the Western traditional metaphysicists and critics of the entire modern paradigm. I wonder if he has ever been exposed to Guénon or related authors? This might be a possible link:

\url{http://www.islambuddhism.com/islam-and-buddhism/}


It is perhaps worth pointing out that DL is only the head of one of the five major schools of Tibet (Gelug, Nyingma, Sakya, Kagyu, Jonang); even though he is respected by most, he is by no means recognised as an infallible authority and even members of his own school may disagree with some of his views.

\hfill

\texttt{Cologero on 2016-12-28 at 20:49 said:}

Several points, O.

The last place I would expect to find an Orthodox theological academy would be in Sweden. Actually, it is not quite Eastern Orthodox, but Oriental ... very ancient. I guess it reflects the changing demographics of Sweden.

Our on-line discussion for the next few months will be dealing with the spiritual impulses from Krisha, Buddha, and Christ. That will give us the opportunity to compare the underlying metaphysical assumptions.

The traditional choices for a spiritual path are the priest and the knight, corresponding to the higher castes. We here, on the other hand, have been describing a third path, more or less outside of the castes --- viz., the Hermetic path. Obviously, that is a very personal choice. I've made my choice and you will make your own.

\hfill

\end{sffamily}
\end{footnotesize}

